% Auxiliary coordinate identities for polynomial PEPS compression.

\subsection{Register identifications and tensor coordinates}

A register has an owner in a set $P$ and a complex inner product space.
For a finite ordered register list $a$, let $M(a)$ be its
right-associated tensor memory, with $M(\varnothing)=\C$ and a
terminal factor $\C$. Write $ab$ for concatenation, and
$J_{a,b}:M(ab)\longrightarrow M(a)\otimes M(b)$ for its canonical
isometry. For a selection $f:P\to\{0,1\}$, let $a_f$ retain the
registers with $f(p)=1$, in their original order, and write
$\bar f=1-f$. The canonical partition isometry is
$U_f^a:M(a)\longrightarrow M(a_f)\otimes M(a_{\bar f})$.
Equal register lists are identified by their canonical isometries,
denoted $E_{a=b}$. All tensor products below retain these conventions.

\begin{theorem}[Equality identifications of tensor memories]
    \label{thm:peps_basis_memory_equalities}
    \lean{TNLean.PEPS.PairEffect.Layout.memCongr_apply_heq,
        TNLean.PEPS.PairEffect.Layout.memCongr_tmul_heq,
        TNLean.PEPS.PairEffect.Layout.memCongr_symm,
        TNLean.PEPS.PairEffect.Layout.appendIso_symm_tmul_heq}
    \leanok
    For equal register lists $a=a'$, the identification $E_{a=a'}$
    preserves each vector after identifying its underlying space, and
    $E_{a'=a}=E_{a=a'}^{-1}$. For $a=a'$ and $b=b'$, the tensor
    $E_{a=a'}x\otimes E_{b=b'}y$ is the same tensor as $x\otimes y$
    under the induced identification. More generally, if $x'$ and $y'$
    agree with $x$ and $y$ under these identifications, then
    $J_{a,b}^{-1}(x\otimes y)$ and
    $J_{a',b'}^{-1}(x'\otimes y')$ agree under the identification
    $M(ab)=M(a'b')$.
\end{theorem}

\begin{proof}
    \leanok
    Identify the equal register lists and their equal vectors.
    Each identification becomes the identity, so all assertions are
    equalities of the same vectors or isometries.
\end{proof}

\begin{theorem}[The final scalar factor]
    \label{thm:peps_basis_final_scalar}
    \lean{TNLean.PEPS.PairEffect.Layout.memCongr_append_nil_appendIso_symm_tmul,
        TNLean.PEPS.PairEffect.Layout.appendIso_symm_tmul_one_eq_memCongr_symm}
    \leanok
    For every register list $a$, vector $x\in M(a)$ and scalar $z\in\C$,
    $E_{a\varnothing=a}J_{a,\varnothing}^{-1}(x\otimes z)=zx$.
    In particular,
    $J_{a,\varnothing}^{-1}(x\otimes1)=E_{a\varnothing=a}^{-1}x$.
\end{theorem}

\begin{proof}
    \leanok
    Induct on $a$, expanding the concatenation isometry on elementary
    tensors. Move the scalar through each tensor factor. Linearity
    gives the identity for arbitrary vectors; set $z=1$ for the
    second assertion.
\end{proof}

\begin{definition}[Restriction after changing owner labels]
    \label{def:peps_basis_owner_restriction}
    \lean{TNLean.PEPS.PairEffect.Layout.restrictMapOwnerIso,
        TNLean.PEPS.PairEffect.Layout.restrict_mapOwner}
    \leanok
    Let $q:P\to Q$ be any map. Write $q_*a$ for the register list
    obtained by replacing each owner $p$ by $q(p)$, retaining the
    ordered register spaces, and let $T_q^a:M(a)\to M(q_*a)$ be the
    tensor product of their identity maps. For a selection $g$ on $Q$,
    $(q_*a)_g=q_*(a_{g\circ q})$. The selected-memory isometry
    $T_{q,g}^a:M(a_{g\circ q})\to M((q_*a)_g)$ is $T_q^{a_{g\circ q}}$
    followed by the identification of these equal register lists.
\end{definition}

\begin{theorem}[Changing owners preserves memory vectors]
    \label{thm:peps_basis_owner_memory}
    \lean{TNLean.PEPS.PairEffect.Layout.mem_mapOwner,
        TNLean.PEPS.PairEffect.Layout.mapOwnerIso_apply_heq}
    \leanok
    For any $q:P\to Q$ and register list $a$, the ordered tensor spaces
    satisfy $M(q_*a)=M(a)$. Under this equality, $T_q^a x=x$ for
    every $x\in M(a)$.
\end{theorem}

\begin{proof}
    \leanok
    Induct on the register list. Relabelling leaves each head space
    unchanged; apply the induction hypothesis to its tail. On the
    resulting tensor product the isometry is the identity, first on
    elementary tensors and then by linearity.
\end{proof}

\begin{theorem}[Partitioning commutes with changing owners]
    \label{thm:peps_basis_owner_partition}
    \lean{TNLean.PEPS.PairEffect.Layout.partitionIso_mapOwner_heq,
        TNLean.PEPS.PairEffect.Layout.partitionIso_mapOwner}
    \leanok
    For every $q:P\to Q$, selection $g$ on $Q$, register list $a$,
    and $x\in M(a)$,
    \begin{align}
        U_g^{q_*a}T_q^a x
         & =(T_{q,g}^a\otimes T_{q,\bar g}^a)U_{g\circ q}^a x.
    \end{align}
    After identifying the selected and complementary tensor spaces,
    this is the same vector as $U_{g\circ q}^a x$.
\end{theorem}

\begin{proof}
    \leanok
    \uses{def:peps_basis_owner_restriction,thm:peps_basis_owner_memory}
    Induct on $a$. The first register has the same selection value
    before and after relabelling. Use the corresponding associativity
    or exchange isometry, and apply the induction hypothesis to the
    tail. The selected-memory isometries preserve the underlying
    vectors, giving the displayed identity.
\end{proof}

\begin{theorem}[A partition with no selected registers]
    \label{thm:peps_basis_empty_selection}
    \lean{TNLean.PEPS.PairEffect.Layout.partitionIso_of_forall_false}
    \leanok
    If $f$ vanishes at the owner of every register in $a$, then
    $a_f=\varnothing$, $a_{\bar f}=a$, and, under these list
    identifications, $U_f^a x=1\otimes x$ for every $x\in M(a)$.
\end{theorem}

\begin{proof}
    \leanok
    The filtered lists have the stated forms. Induct on $a$, using
    the complementary-head partition formula. Exchange the scalar
    unit with each head factor; linearity extends the identity from
    elementary tensors to every vector.
\end{proof}

\begin{theorem}[Grouping equal register lists]
    \label{thm:peps_basis_grouping_equality}
    \lean{TNLean.PEPS.PairEffect.groupByPartyIso_apply_heq}
    \leanok
    Let $\mathbf p$ be an ordered list without repetitions containing
    every owner of registers in $a$ and $b$. Suppose $a=b$ and that
    $x\in M(a)$ and $y\in M(b)$ agree under this identification.
    Write $G_{\mathbf p,a}$ for the isometry that successively
    collects the registers of each party in the prescribed order.
    Then $G_{\mathbf p,a}x$ and $G_{\mathbf p,b}y$ agree under
    the equality of their grouped register lists. The list need not contain parties absent from
    the register memories.
\end{theorem}

\begin{proof}
    \leanok
    Identify $a=b$ and $x=y$. The two isometries give the same successive
    partitions of the same ordered registers.
\end{proof}

\begin{theorem}[The first party of a grouped tensor]
    \label{thm:peps_basis_first_party_grouping}
    \lean{TNLean.PEPS.PairEffect.groupByPartyIso_cons_partition_tmul}
    \leanok
    \uses{def:peps_party_partition}
    Let $P$ be any party set and let $(p,\mathbf p)$ be an ordered
    list without repetitions containing every owner in the ordered
    register list $a$. Write $a_p$ and $a_{\bar p}$ for its registers
    owned by $p$ and its remaining registers, respectively. Let $U_p^a$
    be the partition isometry and $G_{\mathbf p,a}$ the grouping
    isometry in the prescribed order. Each grouped memory has one
    memory factor per listed party, with a terminal factor $\C$.

    Since $p$ is absent from $\mathbf p$, removing its registers
    leaves the registers of each party in $\mathbf p$ unchanged.
    Let $E$ identify the resulting equal grouped tail memories.
    For every $x\in M(a_p)$ and $y\in M(a_{\bar p})$,
    \begin{align}
        G_{(p,\mathbf p),a}(U_p^a)^{-1}(x\otimes y)
         & =x\otimes E G_{\mathbf p,a_{\bar p}}y.
    \end{align}
    The list need not contain parties absent from $a$; its tail
    may be empty. No dimension or nonvanishing condition is imposed.
\end{theorem}

\begin{proof}
    \leanok
    \uses{def:peps_party_partition}
    By the recursive grouping identity,
    $G_{(p,\mathbf p),a}=(1\otimes E G_{\mathbf p,a_{\bar p}})U_p^a$.
    Cancel $U_p^a(U_p^a)^{-1}$ and apply the tensor product map to
    $x\otimes y$.
\end{proof}

\begin{definition}[Joining two partitioned memories]
    \label{def:peps_basis_partition_join}
    \lean{TNLean.PEPS.PairEffect.Layout.appendPartitionIso,
        TNLean.PEPS.PairEffect.Layout.appendPartitionIso_tmul}
    \leanok
    For register lists $a_1,a_2,b_1,b_2$, let $K$ be the isometry
    that exchanges the middle two factors and concatenates the two
    resulting pairs. For every $x_i\in M(a_i)$ and $y_i\in M(b_i)$,
    \begin{align}
        K\bigl((x_1\otimes x_2)\otimes(y_1\otimes y_2)\bigr)
         & =J_{a_1,b_1}^{-1}(x_1\otimes y_1)
        \otimes J_{a_2,b_2}^{-1}(x_2\otimes y_2).
    \end{align}
\end{definition}

\begin{theorem}[Partitioning a concatenation]
    \label{thm:peps_basis_partition_concat}
    \lean{TNLean.PEPS.PairEffect.Layout.partitionIso_append_tmul}
    \leanok
    For any selection $f$, register lists $a,b$, and vectors
    $x\in M(a)$, $y\in M(b)$, identify $(ab)_f=a_fb_f$ and
    $(ab)_{\bar f}=a_{\bar f}b_{\bar f}$. With $K$ from
    Definition~\ref{def:peps_basis_partition_join} for these four lists,
    \begin{align}
        U_f^{ab}J_{a,b}^{-1}(x\otimes y)
         & =K(U_f^a x\otimes U_f^b y).
    \end{align}
\end{theorem}

\begin{proof}
    \leanok
    \uses{def:peps_basis_partition_join}
    Induct on $a$ and extend by linearity in $x$. A selected head
    adjoins its factor to the first joined memory; an excluded head
    adjoins it to the second. In each case, associativity and the
    exchange of the middle factors reduce the assertion to the tail.
\end{proof}

\begin{definition}[Bases adapted to a partition]
    \label{def:peps_basis_partition_bases}
    \lean{TNLean.PEPS.PairEffect.Layout.partitionBasis,
        TNLean.PEPS.PairEffect.Layout.partitionAppendBasis}
    \leanok
    Choose orthonormal bases $b_A,b_E$ of $M(a_f),M(a_{\bar f})$,
    and $b_D,b_F$ of $M(b_f),M(b_{\bar f})$, indexed by finite sets
    $X,Y,I,L$, respectively. Their partition bases on the original
    memories are $D^a_{x,y}=(U_f^a)^{-1}(b_A(x)\otimes b_E(y))$
    and $D^b_{i,j}=(U_f^b)^{-1}(b_D(i)\otimes b_F(j))$.
    On the selected and complementary concatenated memories, take
    $C_{x,i}=J_{a_f,b_f}^{-1}(b_A(x)\otimes b_D(i))$ and
    $F_{y,j}=J_{a_{\bar f},b_{\bar f}}^{-1}(b_E(y)\otimes b_F(j))$,
    transported along the corresponding list equalities.
\end{definition}

\begin{theorem}[Coordinates before and after a joint partition]
    \label{thm:peps_basis_partition_coordinates}
    \lean{TNLean.PEPS.PairEffect.Layout.partitionIso_partitionBasis_tmul,
        TNLean.PEPS.PairEffect.Layout.partitionBasis_coordinates,
        TNLean.PEPS.PairEffect.Layout.partitionBasis_coordinates_of_eq}
    \leanok
    With the bases of Definition~\ref{def:peps_basis_partition_bases},
    every $x\in X$, $y\in Y$, $i\in I$, $j\in L$ satisfies
    \begin{align}
        U_f^{ab}J_{a,b}^{-1}(D^a_{x,y}\otimes D^b_{i,j})
         & =C_{x,i}\otimes F_{y,j},\label{eq:peps_basis_partition_columns} \\
        \langle D^a_{x,y}\otimes D^b_{i,j},J_{a,b}z\rangle
         & =\langle C_{x,i}\otimes F_{y,j},U_f^{ab}z\rangle
        \label{eq:peps_basis_partition_coefficients}
    \end{align}
    The second identity holds for every $z\in M(ab)$. It also holds
    for a register list $c$ with a specified equality $c=ab$ and
    $z\in M(c)$: apply $E_{c=ab}$ before $J_{a,b}$ on the left,
    and transport $C$ and $F$ to $M(c_f)$ and $M(c_{\bar f})$ on
    the right, where the partition is $U_f^c$.
\end{theorem}

\begin{proof}
    \leanok
    \uses{def:peps_basis_partition_bases,thm:peps_basis_partition_concat}
    Apply the concatenation identity to the two partition-basis
    columns and cancel the inverse partition isometries. This gives
    (\ref{eq:peps_basis_partition_columns}). Taking inner products
    and using the isometries yields
    (\ref{eq:peps_basis_partition_coefficients}). Identify $c=ab$
    for the final assertion.
\end{proof}

\begin{definition}[Physical coordinates on the selected output]
    \label{def:peps_basis_partition_output_coordinates}
    \lean{TNLean.PEPS.PairEffect.Layout.partitionOutputIso}
    \leanok
    For any register lists $a,b$ and selection $f$, let $b_A$ be
    an orthonormal basis of
    $M(a_f)$ indexed by a finite set $X$. Denote its coordinate
    isometry by $b_A^\#$. The selected output isometry is
    $O=(b_A^\#\otimes1)J_{a_f,b_f}$, after identifying
    $(ab)_f=a_fb_f$.
\end{definition}

\begin{theorem}[Coordinates of the selected output]
    \label{thm:peps_basis_partition_output_coordinates}
    \lean{TNLean.PEPS.PairEffect.Layout.partitionOutputIso_basis,
        TNLean.PEPS.PairEffect.Layout.partitionOutputIso_coordinates}
    \leanok
    \uses{def:peps_basis_partition_output_coordinates,def:peps_basis_partition_bases}
    With the selected output isometry just defined, for any
    orthonormal basis $b_D$ of $M(b_f)$
    indexed by a finite set $I$, and any $x\in X$, $i\in I$,
    \begin{align}
        O C_{x,i} & =e_x\otimes b_D(i),        \\
        \langle e_x\otimes b_D(i),Oz\rangle
                  & =\langle C_{x,i},z\rangle.
    \end{align}
    Here $e_x$ is the standard basis of $\C^X$,
    $C_{x,i}=J_{a_f,b_f}^{-1}(b_A(x)\otimes b_D(i))$ in the
    identified selected memory, and $z\in M((ab)_f)$ is arbitrary.
\end{theorem}

\begin{proof}
    \leanok
    \uses{def:peps_basis_partition_output_coordinates,def:peps_basis_partition_bases}
    The concatenation isometry sends $C_{x,i}$ to
    $b_A(x)\otimes b_D(i)$. Taking the coordinates of the first
    factor gives $e_x\otimes b_D(i)$. Preservation of inner products
    gives the second identity.
\end{proof}

\begin{theorem}[Tensor coordinates under owner relabelling]
    \label{thm:peps_basis_owner_coordinates}
    \lean{TNLean.PEPS.PairEffect.Layout.mapOwner_tensorBasis_coordinates}
    \leanok
    Let $q:P\to Q$, let $a,b$ be register lists, and let $b_A,b_D$
    be orthonormal bases of $M(q_*a),M(q_*b)$ indexed by finite sets
    $X,I$. For every $z\in M(ab)$, $x\in X$, $i\in I$,
    \begin{align}
        \langle (T_q^a)^{-1}b_A(x)\otimes(T_q^b)^{-1}b_D(i),J_{a,b}z\rangle
         & =\langle b_A(x)\otimes b_D(i),
        J_{q_*a,q_*b}T_q^{ab}z\rangle,
    \end{align}
    with $q_*(ab)=(q_*a)(q_*b)$ identified on the right.
\end{theorem}

\begin{proof}
    \leanok
    Changing owners commutes with concatenation. Apply that identity
    to the two transported basis columns, then use preservation of
    inner products by the owner and concatenation isometries.
\end{proof}

\subsubsection{Physical register lists and regional dimensions}

For a nonnegative integer $d$, write $[d]=\{0,\ldots,d-1\}$.
If $d=0$, this set is empty. Let $\mathbf p=(p_0,\ldots,p_{m-1})$
be an ordered list of parties. Repetitions are permitted unless
excluded explicitly.

\begin{theorem}[Owners and restrictions of physical lists]
    \label{thm:peps_basis_physical_list_restriction}
    \lean{TNLean.PEPS.PairEffect.familyPhysicalLayout_owners,
        TNLean.PEPS.PairEffect.restrict_familyPhysicalLayout}
    \leanok
    \uses{def:peps_family_physical_output}
    Let $d_p$ be arbitrary nonnegative local dimensions, and write
    $a_{\mathbf d}(\mathbf p)$ for the register list with space
    $\C^{d_{p_i}}$ in position $i$. The owner sequence of
    $a_{\mathbf d}(\mathbf p)$ is precisely
    $\mathbf p$. Restriction by any selection $f$ retains exactly
    the corresponding subsequence:
    \begin{align}
        (a_{\mathbf d}(\mathbf p))_f & =a_{\mathbf d}(\mathbf p_f).
    \end{align}
\end{theorem}

\begin{proof}
    \leanok
    Taking owners reads the same ordered party list. Filtering the
    registers commutes with assigning their spaces, giving the
    restriction identity.
\end{proof}

\begin{definition}[Physical coordinates of a region]
    \label{def:peps_basis_region_coordinates}
    \uses{def:peps_family_physical_output,thm:peps_basis_physical_list_restriction}
    \lean{TNLean.PEPS.PairEffect.filteredPartyEquiv,
        TNLean.PEPS.PairEffect.familyRestrictedPhysicalBasis}
    \leanok
    Suppose $\mathbf p$ has no repetitions and contains every party
    of $P$, and let $A\subseteq P$ be finite. Reading the party at
    each position of the filtered list $\mathbf p_A$ is a bijection
    from its positions to $A$. Reindex the ordered standard tensor
    basis of these registers by this bijection. Its indices are
    assignments $x_p\in[d_p]$ for $p\in A$.
    Transport along the physical-list restriction identity to obtain
    bases on the actual memories selected by $A$. The common-dimension
    case is obtained by setting $d_p=d$.
\end{definition}

\begin{theorem}[Physical dimensions of lists and regions]
    \label{thm:peps_basis_physical_dimensions}
    \lean{TNLean.PEPS.PairEffect.length_filter_region,
        TNLean.PEPS.PairEffect.finrank_restrict_familyPhysicalLayout,
        TNLean.PEPS.PairEffect.card_familyRegionalPhysicalIndex,
        TNLean.PEPS.PairEffect.finrank_restrict_familyPhysicalLayout_le}
    \leanok
    For arbitrary nonnegative local dimensions $d_p$ and finite $A\subseteq P$,
    \begin{align}
        \#\prod_{p\in A}[d_p] & =\prod_{p\in A}d_p.
    \end{align}
    If $\mathbf p$ has no repetitions and contains every party of
    $P$, then $|\mathbf p_A|=|A|$ and
    \begin{align}
        \dim M((a_{\mathbf d}(\mathbf p))_A) & =\prod_{p\in A}d_p.
    \end{align}
    In particular, for a common dimension $d_p=d$, both the number
    of regional configurations and the selected memory dimension are $d^{|A|}$.
    If $D$ is a nonnegative integer and $d_p\le D$ for every
    $p\in A$, the selected memory dimension is at most $D^{|A|}$.
    Empty products have value $1$,
    including when $A$ is empty. If $d_p=0$ for some $p\in A$,
    the regional dimension is zero.
\end{theorem}

\begin{proof}
    \leanok
    \uses{def:peps_family_physical_output,
        thm:peps_basis_physical_list_restriction,def:peps_basis_region_coordinates}
    Count the assignments indexing each tensor basis. An exhaustive
    ordering without repetitions contains each element of $A$ once,
    so its filtered list has length $|A|$. Dimension equals the
    number of basis elements. Multiply the inequalities $d_p\le D$
    over $A$ to obtain the bound.
\end{proof}

\begin{definition}[Physical and regional coordinate isometries]
    \label{def:peps_basis_physical_readouts}
    \lean{TNLean.PEPS.PairEffect.familyLabelledPhysicalReadout,
        TNLean.PEPS.PairEffect.familyRegionalPhysicalIndexEquiv,
        TNLean.PEPS.PairEffect.familyRegionalPhysicalReadout}
    \leanok
    For any set $P$ and finite $A\subseteq P$, splitting an assignment
    into its restrictions to $A$ and $P\setminus A$ is a bijection,
    for both constant
    dimension $d$ and arbitrary local dimensions $d_p$.
    Now suppose $\mathbf p$ lists $P$ without repetitions, and let
    $b$ be a private register list with finite-dimensional memory
    $M(b)$ of dimension $r$. Fix its standard orthonormal basis
    $\gamma_0,\ldots,\gamma_{r-1}$. Let $E_x$ denote the labelled
    physical standard basis, and define the coordinate isometry
    $R:M(a b)\to\C^{\mathcal X\times[r]}$ by
    \begin{align}
        R(v)(x,i)=\langle E_x\otimes\gamma_i,J_{a,b}v\rangle,
    \end{align}
    where $a$ is the chosen physical list and $\mathcal X$ its
    assignment set. The regional coordinate isometry
    $R_A:M(ab)\to\C^{\mathcal X_A\times(\mathcal X_{P\setminus A}\times[r])}$
    is obtained by separating the physical assignment into its
    regional and exterior restrictions, with the private index
    retained in the latter coordinate pair.
\end{definition}

\begin{theorem}[Regional coordinate formulas and contraction bounds]
    \label{thm:peps_basis_regional_readouts}
    \lean{TNLean.PEPS.PairEffect.norm_familyRegionalPhysicalReadout_le_one,
        TNLean.PEPS.PairEffect.familyRegionalPhysicalReadout_apply}
    \leanok
    Under the hypotheses of
    Definition~\ref{def:peps_basis_physical_readouts}, let $s,t$ be
    assignments on $A,P\setminus A$, respectively, and let $s\sqcup t$
    be their unique common assignment on $P$. For every
    $v\in M(ab)$ and $i\in[r]$, both choices of physical dimensions
    satisfy $R_A(v)(s,(t,i))=R(v)(s\sqcup t,i)$ and
    $\|R_A\|_{\rm op}\le1$.
\end{theorem}

\begin{proof}
    \leanok
    \uses{def:peps_basis_physical_readouts}
    Evaluate the reindexed product-basis coordinates at
    $(s,(t,i))$. This is the original product-basis coordinate at
    $(s\sqcup t,i)$. The coordinate map is an isometry, which gives
    the operator-norm bound, also for a zero-dimensional domain.
\end{proof}

\subsubsection{Ordered party tensor coordinates}

Let $a_p$ select all registers of $a$ owned by $p$. For every $p$,
choose an orthonormal basis $b_p$ of $M(a_p)$ indexed by a finite
set $I_p$. For any ordered party list
$\mathbf p=(p_0,\ldots,p_{m-1})$, let $B^{\mathbf p}$ be its
right-associated product basis on
$M(a_{p_0})\otimes\cdots\otimes M(a_{p_{m-1}})\otimes\C$.

\begin{theorem}[Recursion for party tensor coordinates]
    \label{thm:peps_basis_party_recursion}
    \lean{TNLean.PEPS.PairEffect.partyListBasis_nil_apply,
        TNLean.PEPS.PairEffect.partyListBasis_repr_nil,
        TNLean.PEPS.PairEffect.partyListBasis_cons_apply,
        TNLean.PEPS.PairEffect.partyListBasis_repr_cons_tmul}
    \leanok
    The empty list has one basis index, with column $1$, and the
    coordinate of $z\in\C$ at this index is $z$. If a party $p$
    is adjoined to the head of an arbitrary list $\mathbf p$, its
    basis column with index $(x_0,x_{\rm tail})$ is
    $b_p(x_0)\otimes B^{\mathbf p}_{x_{\rm tail}}$.
    For every $u\in M(a_p)$ and tail vector $v$,
    \begin{align}
        \langle B^{p\mathbf p}_{(x_0,x_{\rm tail})},u\otimes v\rangle
         & =\langle b_p(x_0),u\rangle
        \langle B^{\mathbf p}_{x_{\rm tail}},v\rangle.
    \end{align}
    The list may have repetitions and may omit parties.
\end{theorem}

\begin{proof}
    \leanok
    The empty basis is the scalar singleton basis. The recursive
    tensor basis has the stated columns. Its coordinate on an
    elementary tensor is the product of the two factor coordinates.
\end{proof}

\begin{theorem}[Matrix entries of ordered local tensor maps]
    \label{thm:peps_basis_party_map_entries}
    \lean{TNLean.PEPS.PairEffect.partyListBasis_repr_tensorPartyMaps,
        TNLean.PEPS.PairEffect.partyLabelledBasis_apply,
        TNLean.PEPS.PairEffect.partyLabelledBasis_repr_tensorPartyMaps}
    \leanok
    Let $a,c$ be register lists, and for every $p$ choose finite
    orthonormal bases $b_p,c_p$ of $M(a_p),M(c_p)$ and a continuous
    linear map $A_p:M(a_p)\to M(c_p)$. For any ordered party list,
    the product map $A_{p_0}\otimes\cdots\otimes A_{p_{m-1}}\otimes1$
    has matrix entries
    \begin{align}
        \left\langle C^{\mathbf p}_x,
        (\bigotimes_{i=0}^{m-1}A_{p_i})B^{\mathbf p}_y\right\rangle
         & =\prod_{i=0}^{m-1}\langle c_{p_i}(x_i),A_{p_i}b_{p_i}(y_i)\rangle.
    \end{align}
    If $P$ is finite and $\mathbf p$ lists all of $P$ without
    repetitions, reindex the columns by party assignments: the
    labelled column at $y$ is the ordered column at
    $(y(p_0),\ldots,y(p_{m-1}))$. In these labelled bases the matrix
    entry at $(x,y)$ equals
    $\prod_{p\in P}\langle c_p(x_p),A_pb_p(y_p)\rangle$.
    No contraction bound on the local maps is imposed.
\end{theorem}

\begin{proof}
    \leanok
    \uses{thm:peps_basis_party_recursion}
    Induct on the party list. Apply the first local map to the first
    basis column and use the product coordinate formula to separate
    it from the tail. The empty product is $1$. For labelled
    coordinates, reindex both bases by the exhaustive ordering and
    replace the product over its positions by the product over $P$.
\end{proof}

These are the tensor-coordinate identities used in
\cite[proof of Theorem~5.2]{OpenAI2026PolynomialPEPS},
\texttt{04-compression.tex}, lines 137--151, 233--267, 383--450 and
565--588.
